% !TeX root = ../papers/nc-p3.tex
\subsection{有限な四辺での同型不変量と最小周期}
\label{sec:curvilinear-window-orbits}
基礎体は標数零の代数閉体とする。定理\ref{thm:finite-square-chart}の三つの
四次元空間を使い、局所環が全て曲線状の場合の同型を数値の一致に帰着する。
重複点の長さと個数には制限を置かない。

\begin{theorem}\label{thm:curvilinear-window-orbits}
四辺が有限な標準AS窓で、点環が
$D=\prod_{a=1}^h k[t_a]/(t_a^{m_a})$、$\sum m_a=20$とする。
以下の手順で得る、長さ順の各標識付けに対する記録の有限集合は
添字保持同型の完全不変量である。記録は長さ、三組の主元行、
零でない座標の位置、非零な体の元の列からなる。
同じ長さの局所因子の置換以外の軌道探索を必要としない。
\end{theorem}
\begin{proof}
局所因子を順に並べ、各因子では $1,t,\ldots,t^{m-1}$ の順を使う。
三位置 $s=0,1,2$ の空間を列行列 $B_s$ で表示する。
最初から行を読んで独立になるたびに採用した四つの行を
$P=(p_1,\ldots,p_4)$ とし、右から主元小行列の逆を掛けて $B_P=I$ とする。
下三角可逆行列による左作用は先頭の全ての区間の階数を保つので $P$ を変えない。

各位置で独立な単元倍と、三位置共通の代数自己同型を作用させる。
因子の標識を固定すると、この群は対角トーラスと下単三角群 $N$ との積である。
$N$ のLie環の生成ベクトルは、各因子の長さを $m$ として
\[
 \begin{array}{ll}
 t^k\text{ の乗法}&(1\le k<m,\ \text{一つの位置に作用}),\\
 t^{k+1}\partial_t &(1\le k<m-1,\ \text{三位置に同時に作用})
 \end{array}
\]
で張られる。実際、定数項1の単元と一次係数1の座標変換の接空間が
この二種類であり、各群は係数を座標とする連結な単冪群である。
トーラスの $(s,a,d)$ 行での重みは $c_{s,a} b_a^d$ である。
ここで $c_{s,a}$ は単元の定数、$b_a$ は $t_a$ の一次係数である。
$m_a=1$ のときは $b_a$ を省いてよい。

各位置内で非主元行 $r$ を昇順に、列 $j$ を降順に取り、$p_j<r$ の座標だけを読む。
$g\in N$ の作用は $B\mapsto (gB)((gB)_P)^{-1}$ であり、
各座標に先行座標の多項式を加える。実際、主元小行列は下単三角で、
その逆の第 $i$ 行の非対角成分は $p_i$ より前の行だけで決まる。
$r$ 行の積で使うのは $p_i<r$ の列に限られ、$j$ より大きい列は先に読んでいる。
従って既読の値を固定する部分群は次の一座標に平行移動で作用する。

現在の安定化Lie環の元 $X$ の平行移動速度は、対応する成分を同じ記号で書けば
\[
                 (XB-B(XB)_P)_{rj}
\]
である。これは線型汎関数で、零なら値を残す。非零なら速度 $v\ne0$ の
$X$ を選び、三位置へ $\exp(-B_{rj}X/v)$ を作用させ、この座標を零にする。
次は速度の核を使う。単冪群では閉部分群も指数・対数でそのLie環と対応する。
ここでは直接、部分群の元 $g$ が既読値を保てば $g^n$ も全て保ち、
$\exp(z\log g)$ の正規化後の各成分が $z$ の多項式なので全ての $z$ で保つ。
逆向きは既読座標への三角な微分方程式の一意性から従う。
従って速度が零の場合にはその値は軌道上で強制される。
帰納的に同一軌道の代表をただ一つ選んでいるので、速度方向の選択は最終値を変えない。
非零速度を使う回数は当初のLie環の次元以下である。

この正規化はトーラスと可換である。トーラスは $N$ を正規化し、
速度が零かどうかと零にした座標を保つので、各段階の代表の一意性を適用できる。
残った非零座標を $z_1,\ldots,z_l$ とする。行 $(a,d)$、主元行 $(a',e)$ の
座標の重みは $c_{s,a}b_a^d/(c_{s,a'}b_{a'}^e)$ である。
これらの重みを行にした整数行列を $W$ とし、
$\ker_{\mathbb Z}(W^{\mathsf t})$ の整数基底 $u$ ごとに $z^u$ を記録する。
主元と台が等しいときこの値の一致はトーラス同値と同値である。
実際、一つのトーラス変数について行を整数の可逆基本変形し、その指数を
一本の行に集める。対応する単項式座標変換も可逆である。
その変数は非零べき乗を通じこの一座標だけに全射に作用するので、
残る行に同じ手順を施せばよい。除いた座標は逆順に合わせられる。
有理核の独立なベクトルを個別に整数化するだけでは整数核の基底を保証しない。
最後に代数同型は同じ長さの局所因子だけを交換できるから、
標識付けの有限集合を比較し、定理\ref{thm:finite-square-chart}を適用する。
\end{proof}

\begin{corollary}\label{cor:curvilinear-window-period}
この枝で最小周期を求めるには、開始位置 $i=0,1,2$ の各々で、実際の辺の
二つの射影を通じて三位置を同じスキームへ移し、その三空間の記録集合を
$K_i$ とする。$K_0=K_1$ なら1、そうでなく $K_0=K_2$ なら2、残りは4である。
\end{corollary}
\begin{proof}
$K_i=K_0$ と $i$ 移動への添字保持同型は定理から同値である。
同型を許す移動量は $\mathbb Z$ の部分群であり、標準四周期性により4を含む。
境界を越える三位置にも辺の射影を用いるので、四移動の点スキーム上の作用を
恒等写像と仮定する必要はない。
\end{proof}

\begin{proposition}\label{prop:finite-point-aut-solvable}
一般の有限四辺でも、各局所点環 $C$ の自己同型群の単位成分は可解である。
従って一つの点環 $D$ の同型型とその表示を固定すれば、三空間の同型判定は
有限個の成分代表と、三角行列に対する上の消去法に帰着する。
\end{proposition}
\begin{proof}
辺は滑らかな六次元多様体内の有限完全交叉で、第一射影は閉埋込みである。
従って $C$ の埋込み次元 $e$ は高々3である。六つの局所方程式の独立な
一次部を形式座標に採用して消去すると、極小表示は
$C=k[[x_1,\ldots,x_e]]/I$、$I\subset (x)^2$ であり、
極小関係空間 $I/(x)I$ の次元は $e$ である（高さも $e$）。

自己同型が余接空間上で恒等なら、極大イデアルの冪による全ての次数商上で
恒等なので作用は単冪である。単位成分が可解でなければ標数零のLevi分解と
半単純Lie環の根部分環から、余接空間に非自明に作用する
$\mathfrak{sl}_2$ が取れる。完全可約性で余接空間を極大イデアル内に
同変に持ち上げ、上の形式変数をその基底に取る。
$W=I/(x)I$ にも同変な関係式代表を取れる。実際、
$(x)^n\subset (x)I$ となる $n$ で切り、$I$ 内の次数 $n$ 以下の多項式から
$W$ への全射を分裂させればよい。線型変数への作用はこの有限次元空間を保つ。
この $e$ 本は中山の補題で $I$ を生成する。

$e\le3$ で非自明な $\mathfrak{sl}_2$ 表現は、二次元標準表現、これと
自明な一次元の和、または最高重み2の三次元既約表現である。
変数の最高重みを $M>0$ とし、最高重みの変数一つだけを残す。
全ての関係がこの直線上で消えれば商が有限にならない。
一次式がないので $W$ に重み $2M$ 以上が必要である。
$e=2$ または $M=2$ なら次元 $e$ の表現にその重みはない。
残る $e=3,M=1$ なら $W$ は三次元既約で、自明な変数の軸への制限は
自明表現への準同型なので零である。やはり有限性に反する。

$D$ に対し、三組の単元群と代数自己同型群の半直積の単位成分は以上により可解である。
各位置の $D$ 上でLie--Kolchinの定理による不変完全旗を取り、極大トーラスの
固有ベクトルで旗に適合する基底を選ぶ。連結可解群は極大トーラスと単冪根基の積で、
後者はこの基底で下単三角となる（旗の順を反転する）。
従って単冪根基のLie環と対角重みを用いれば先の消去証明がそのまま働く。
有限成分群の各代表を先に作用させれば全群の軌道が得られる。
この還元で非曲線状の局所環の同型型そのものを列挙したことにはならない。
\end{proof}

\path{computations/verify_curvilinear_window_orbits.py} は被約二十点の対称AS例と、
局所長さ4の点を含む長さ20の三空間で、基底変更・単元倍・共通の座標置換による
値の不変性を検査する。後者で一位置だけの接係数を変える対照では
標識付きの値が異なる。後者は軌道計算の試験入力であり、AS性は用いない。
